% Generated by Naberius. Content is escaped; shell escape is disabled.
\documentclass[12pt,a4paper]{article}
\usepackage{fontspec}
\defaultfontfeatures{Ligatures=TeX}
\setmainfont[Path=fonts/,BoldFont=lmroman10-bold.otf,ItalicFont=lmroman10-italic.otf,BoldItalicFont=lmroman10-bolditalic.otf]{lmroman10-regular.otf}
\usepackage[spanish]{babel}
\usepackage[margin=25.0mm,headheight=16pt]{geometry}
\usepackage{graphicx,amsmath,amssymb,tabularx,booktabs}
\usepackage{fancyhdr}
\usepackage[unicode,hidelinks]{hyperref}
\hypersetup{pdftitle={El oficio de escribir},pdfauthor={Documento de demostración},pdfkeywords={}}
\addto\captionsspanish{\renewcommand{\figurename}{Figura}\renewcommand{\tablename}{Tabla}\renewcommand{\contentsname}{Índice}}
\setcounter{secnumdepth}{4}
\pagestyle{fancy}
\fancyhf{}\fancyhead[L]{\small Prueba de \LaTeX{}}\fancyfoot[C]{\small Naberius\quad \thepage}\renewcommand{\headrulewidth}{0.4pt}
\fancypagestyle{plain}{\fancyhf{}\fancyhead[L]{\small Prueba de \LaTeX{}}\fancyfoot[C]{\small Naberius\quad \thepage}\renewcommand{\headrulewidth}{0.4pt}}
\linespread{1.15}
\setlength{\parindent}{0.0mm}
\setlength{\parskip}{7.0pt}
\begin{document}
\begin{abstract}Prueba tipográfica de \LaTeX{}: español, áéíóú y referencias.\end{abstract}
\fontsize{12.0}{14.399999999999999}\selectfont
{\LARGE\bfseries El oficio de escribir\par}

{\large Del primer borrador a una idea que se entiende.\par}

\section{Una idea necesita espacio}\label{b620E1581-D159-4E45-AFE1-A00A37AA7FDE}

Escribir es dar forma a lo que pensamos. Este documento de demostración muestra cómo organizar un texto, enlazar sus fuentes y cuidar su presentación sin interrumpir la escritura.

\subsection{Estructura antes que apariencia}\label{b1C27E752-BF49-42A6-BAA8-2B6CAB0DCA85}

Usa los estilos de la barra para expresar qué significa cada párrafo. El esquema se construye a partir de tus títulos; la plantilla se ocupa de la tipografía.

Inserta una tabla, una figura o una ecuación desde el menú Insertar.

\begin{itemize}
\item Guarda tu trabajo como .nab: contenido, imágenes y referencias viajan juntos.

\end{itemize}
\section{Un documento, dos maneras de verlo}\label{b2CE21FB0-D229-480C-8733-6693020FC6F3}

Alterna entre páginas y lienzo continuo. El texto y el historial son los mismos; la paginación definitiva corresponde al PDF compuesto con \LaTeX{}.

\LaTeX{} y \TeX{} conservan su composición tipográfica.

\appendix
\section{Datos complementarios}\label{bDC535091-139F-4D32-A531-DB1F9D43CE40}

\par\medskip\noindent\rule{\linewidth}{0.4pt}\par\medskip

\begin{table}[htbp]\centering\caption{Resultados de prueba}\label{n422032C5-D2F7-42AB-A198-D16BE70AED0F}
\begin{tabularx}{\linewidth}{|>{\raggedright\arraybackslash}X|>{\raggedright\arraybackslash}X|}\hline
\textbf{Variable} & \textbf{Valor} \\ \hline
 &  \\ \hline
\end{tabularx}\par{\small }\end{table}

\begin{equation}\sqrt{\alpha  + \beta }\label{n1F54088C-8A93-4396-B26F-2899DD6DFA3A}\end{equation}

\end{document}